\documentclass[12pt,table,xcolor={dvipsnames}]{beamer}
\usetheme{Pittsburgh}
\usecolortheme{seagull}
%\usepackage[utf8]{inputenc}
\usepackage{fontspec}
\usepackage[brazilian]{babel}
\usepackage{amsmath}
\usepackage{listings}
\usepackage{multirow}
\usepackage{amsfonts}
\usepackage{amssymb}
\usepackage{lstlinebgrd}
\usepackage{graphicx}
\author{Listas Utilizando Vetores}
\title{Estruturas de Dados}
%\setbeamercovered{transparent}
\setbeamertemplate{navigation symbols}{}
%\logo{\includegraphics[scale=0.015]{Brasao_UFSC.png}\includegraphics[scale=0.2]{brasao_PPGCC.jpg}}
\institute{Departamento de Informática e de Estatística \\ Prof. Jean Everson Martina \\ Prof. Aldo von Wangenheim}
\date{2016.2}
\subject{}
\usebackgroundtemplate{\includegraphics[width=\paperwidth,
height=\paperheight]{../reusable_images/fundo_UFSC.png}}
\begin{document}

{
\usebackgroundtemplate{\includegraphics[width=\paperwidth,
height=\paperheight]{../reusable_images/fundo_capa.png}}
\begin{frame}
\titlepage
\includegraphics[scale=0.3]{../reusable_images/brasao_INE.png}
\end{frame}
}

\begin{frame}{Conceito de Listas}
\begin{block}{Lista}
Uma Lista é um conjunto de dados dispostos e/ou acessáveis em uma sequência determinada.
\end{block}
\end{frame}

\begin{frame}[fragile]{Definições de Lista}
\begin{itemize}
\item Lista é um conjunto de dados que pode possuir uma ordem intrínseca (Lista Ordenada) ou não;
\item Este conjunto de dados pode ocupar espaços de memória fisicamente consecutivos, espelhando a sua ordem, ou não;
\item Se os dados estiverem dispersos fisicamente, para que este conjunto seja uma lista, ele deve possuir operações e informações adicionais que permitam que seja tratado como tal (Lista Encadeada).
\end{itemize}
\end{frame}

\begin{frame}[fragile]{Definições de Lista}
\begin{block}{Lista}
A Lista é uma estrutura de dados cujo funcionamento é inspirado no de uma lista "natural".
\end{block}
\begin{itemize}
\item Uma lista pode ser ordenada ou não:
\item Quando for ordenada, pode o ser por alguma característica intrínseca dos dados (ex: ordem alfabética);
\item Ela pode também refletir a ordem cronológica (ordem de inserção) dos dados.
\end{itemize}
\end{frame}

\begin{frame}[fragile]{Listas Usando Vetores}
\begin{columns}
\column{.5\textwidth}
\begin{itemize}
\item Vetores possuem um espaço limitado para armazenar dados;
\item Necessitamos definir um espaço grande o suficiente para a nossa lista;
\item Necessitamos de um indicador de qual elemento do vetor é o atual último da lista.
\end{itemize}
\column{.5\textwidth}
\begin{center}
\begin{tabular}{| p{.5cm} |p{.5cm}c }
\cline{1-1}
 \cellcolor{SpringGreen}  & &\\ \cline{1-1}
\cellcolor{GreenYellow}  & &\\ \cline{1-1}
\cellcolor{SpringGreen}  & &\\ \cline{1-1}
\cellcolor{GreenYellow}   & &\\ \cline{1-1}
\cellcolor{SpringGreen}  & &\\ \cline{1-2}
\cellcolor{GreenYellow}  & \multicolumn{1}{|p{.5cm}|} {int} & \\ \cline{1-2}
\end{tabular}
\end{center}
\end{columns}
\end{frame}

\begin{frame}[fragile]{Listas Usando Vetores}
\begin{columns}
\column{.5\textwidth}
\begin{itemize}
\item Vetores possuem um espaço limitado para armazenar dados;
\item Necessitamos definir um espaço grande o suficiente para a nossa lista;
\item Necessitamos de um indicador de qual elemento do vetor é o atual último da lista.
\item {\color{red} Lista Vazia!}
\end{itemize}
\column{.5\textwidth}
\begin{center}
\begin{tabular}{| p{.5cm} |p{.5cm}c }
\cline{1-1}
 \cellcolor{SpringGreen}  & &\\ \cline{1-1}
\cellcolor{GreenYellow}  & &\\ \cline{1-1}
\cellcolor{SpringGreen}  & &\\ \cline{1-1}
\cellcolor{GreenYellow}   & &\\ \cline{1-1}
\cellcolor{SpringGreen}  & &\\ \cline{1-2}
\cellcolor{GreenYellow}  & \multicolumn{1}{|p{.5cm}|} {-1} & \\ \cline{1-2}
\end{tabular}
\end{center}
\end{columns}
\end{frame}

\begin{frame}[fragile]{Listas Usando Vetores}
\begin{columns}
\column{.5\textwidth}
\begin{itemize}
\item Vetores possuem um espaço limitado para armazenar dados;
\item Necessitamos definir um espaço grande o suficiente para a nossa lista;
\item Necessitamos de um indicador de qual elemento do vetor é o atual último da lista.
\end{itemize}
\column{.5\textwidth}
\begin{center}
\begin{tabular}{| p{.5cm} |p{.5cm}c }
\cline{1-1}
 \cellcolor{SpringGreen}  & &\\ \cline{1-1}
\cellcolor{GreenYellow}  & &\\ \cline{1-1}
\cellcolor{SpringGreen}  & &\\ \cline{1-1}
\cellcolor{GreenYellow}   & &\\ \cline{1-1}
\cellcolor{SpringGreen}  & &\\ \cline{1-2}
\cellcolor{GreenYellow} 20 & \multicolumn{1}{|p{.5cm}|} {0} & \\ \cline{1-2}
\end{tabular}
\end{center}
\end{columns}
\end{frame}

\begin{frame}[fragile]{Listas Usando Vetores}
\begin{columns}
\column{.5\textwidth}
\begin{itemize}
\item Vetores possuem um espaço limitado para armazenar dados;
\item Necessitamos definir um espaço grande o suficiente para a nossa lista;
\item Necessitamos de um indicador de qual elemento do vetor é o atual último da lista.
\end{itemize}
\column{.5\textwidth}
\begin{center}
\begin{tabular}{| p{.5cm} |p{.5cm}c }
\cline{1-1}
 \cellcolor{SpringGreen}  & &\\ \cline{1-1}
\cellcolor{GreenYellow}  & &\\ \cline{1-1}
\cellcolor{SpringGreen}  & &\\ \cline{1-1}
\cellcolor{GreenYellow}   & &\\ \cline{1-1}
\cellcolor{SpringGreen} 55 & &\\ \cline{1-2}
\cellcolor{GreenYellow} 20 & \multicolumn{1}{|p{.5cm}|} {1} & \\ \cline{1-2}
\end{tabular}
\end{center}
\end{columns}
\end{frame}

\begin{frame}[fragile]{Listas Usando Vetores}
\begin{columns}
\column{.5\textwidth}
\begin{itemize}
\item Vetores possuem um espaço limitado para armazenar dados;
\item Necessitamos definir um espaço grande o suficiente para a nossa lista;
\item Necessitamos de um indicador de qual elemento do vetor é o atual último da lista.
\end{itemize}
\column{.5\textwidth}
\begin{center}
\begin{tabular}{| p{.5cm} |p{.5cm}c }
\cline{1-1}
 \cellcolor{SpringGreen}  & &\\ \cline{1-1}
\cellcolor{GreenYellow}  & &\\ \cline{1-1}
\cellcolor{SpringGreen}  & &\\ \cline{1-1}
\cellcolor{GreenYellow}  4 & &\\ \cline{1-1}
\cellcolor{SpringGreen} 55 & &\\ \cline{1-2}
\cellcolor{GreenYellow} 20 & \multicolumn{1}{|p{.5cm}|} {2} & \\ \cline{1-2}
\end{tabular}
\end{center}
\end{columns}
\end{frame}

\begin{frame}[fragile]{Listas Usando Vetores}
\begin{columns}
\column{.5\textwidth}
\begin{itemize}
\item Vetores possuem um espaço limitado para armazenar dados;
\item Necessitamos definir um espaço grande o suficiente para a nossa lista;
\item Necessitamos de um indicador de qual elemento do vetor é o atual último da lista.
\end{itemize}
\column{.5\textwidth}
\begin{center}
\begin{tabular}{| p{.5cm} |p{.5cm}c }
\cline{1-1}
 \cellcolor{SpringGreen}  & &\\ \cline{1-1}
\cellcolor{GreenYellow}  & &\\ \cline{1-1}
\cellcolor{SpringGreen}  12& &\\ \cline{1-1}
\cellcolor{GreenYellow}  4 & &\\ \cline{1-1}
\cellcolor{SpringGreen} 55 & &\\ \cline{1-2}
\cellcolor{GreenYellow} 20 & \multicolumn{1}{|p{.5cm}|} {3} & \\ \cline{1-2}
\end{tabular}
\end{center}
\end{columns}
\end{frame}

\begin{frame}[fragile]{Listas Usando Vetores}
\begin{columns}
\column{.5\textwidth}
\begin{itemize}
\item Vetores possuem um espaço limitado para armazenar dados;
\item Necessitamos definir um espaço grande o suficiente para a nossa lista;
\item Necessitamos de um indicador de qual elemento do vetor é o atual último da lista.
\end{itemize}
\column{.5\textwidth}
\begin{center}
\begin{tabular}{| p{.5cm} |p{.5cm}c }
\cline{1-1}
 \cellcolor{SpringGreen}  & &\\ \cline{1-1}
\cellcolor{GreenYellow}  89& &\\ \cline{1-1}
\cellcolor{SpringGreen}  12& &\\ \cline{1-1}
\cellcolor{GreenYellow}  4 & &\\ \cline{1-1}
\cellcolor{SpringGreen} 55 & &\\ \cline{1-2}
\cellcolor{GreenYellow} 20 & \multicolumn{1}{|p{.5cm}|} {4} & \\ \cline{1-2}
\end{tabular}
\end{center}
\end{columns}
\end{frame}

\begin{frame}[fragile]{Listas Usando Vetores}
\begin{columns}
\column{.5\textwidth}
\begin{itemize}
\item Vetores possuem um espaço limitado para armazenar dados;
\item Necessitamos definir um espaço grande o suficiente para a nossa lista;
\item Necessitamos de um indicador de qual elemento do vetor é o atual último da lista.
\end{itemize}
\column{.5\textwidth}
\begin{center}
\begin{tabular}{| p{.5cm} |p{.5cm}c }
\cline{1-1}
 \cellcolor{SpringGreen} 24 & &\\ \cline{1-1}
\cellcolor{GreenYellow}  89& &\\ \cline{1-1}
\cellcolor{SpringGreen}  12& &\\ \cline{1-1}
\cellcolor{GreenYellow}  4 & &\\ \cline{1-1}
\cellcolor{SpringGreen} 55 & &\\ \cline{1-2}
\cellcolor{GreenYellow} 20 & \multicolumn{1}{|p{.5cm}|} {5} & \\ \cline{1-2}
\end{tabular}
\end{center}
\end{columns}
\end{frame}

\begin{frame}[fragile]{Listas Usando Vetores}
\begin{columns}
\column{.5\textwidth}
\begin{itemize}
\item Vetores possuem um espaço limitado para armazenar dados;
\item Necessitamos definir um espaço grande o suficiente para a nossa lista;
\item Necessitamos de um indicador de qual elemento do vetor é o atual último da lista.
\item {\color{red} Lista Cheia!}
\end{itemize}
\column{.5\textwidth}
\begin{center}
\begin{tabular}{| p{.5cm} |p{.5cm}c }
\cline{1-1}
 \cellcolor{SpringGreen} 24 & &\\ \cline{1-1}
\cellcolor{GreenYellow}  89& &\\ \cline{1-1}
\cellcolor{SpringGreen}  12& &\\ \cline{1-1}
\cellcolor{GreenYellow}  4 & &\\ \cline{1-1}
\cellcolor{SpringGreen} 55 & &\\ \cline{1-2}
\cellcolor{GreenYellow} 20 & \multicolumn{1}{|p{.5cm}|} {5} & \\ \cline{1-2}
\end{tabular}
\end{center}
\end{columns}
\end{frame}

\begin{frame}[fragile]{Modelagem da Lista de Vetor}
\lstset{language=C++,
          keywordstyle=\color{blue}\ttfamily,
          stringstyle=\color{red}\ttfamily,
          commentstyle=\color{green}\ttfamily,
          breaklines=true,
          basicstyle=\ttfamily\footnotesize
          }

\begin{itemize}
\item Aspecto Estrutural:
\begin{itemize}
\item Necessitamos de um vetor para armazenar as informações;
\item Necessitamos de um indicador da posição atual do ultimo elemento da lista;
\end{itemize}
\begin{lstlisting}
constantes MAXLISTA = 100;

classe Lista {
 T _dados[MAXLISTA]; // Vetor Estatico
 inteiro _ultimo;
};
\end{lstlisting}
\end{itemize}
\end{frame}

\begin{frame}[fragile]{Modelagem da Lista de Vetor}
\lstset{language=C++,
          keywordstyle=\color{blue}\ttfamily,
          stringstyle=\color{red}\ttfamily,
          commentstyle=\color{green}\ttfamily,
          breaklines=true,
          basicstyle=\ttfamily\footnotesize
          }

\begin{itemize}
\item Aspecto Estrutural:
\begin{itemize}
\item Necessitamos de um vetor para armazenar as informações;
\item Necessitamos de um indicador da posição atual do ultimo elemento da lista;
\end{itemize}
\begin{lstlisting}
constantes MAXLISTA = 100;

classe Lista {
 T* _dados; // _dados = new T[MAXLISTA];
 inteiro _ultimo;
 inteiro _tam = MAXLISTA;
};
\end{lstlisting}
\end{itemize}
\end{frame}

\begin{frame}[fragile]{Modelagem da Lista de Vetor}
\lstset{language=C++,
          keywordstyle=\color{blue}\ttfamily,
          stringstyle=\color{red}\ttfamily,
          commentstyle=\color{green}\ttfamily,
          breaklines=true,
          basicstyle=\ttfamily\footnotesize
          }

\begin{itemize}
\item Aspecto Funcional:
\begin{itemize}
\item Temos que colocar e retirar dados da lista;
\item Temos que testar se a lista está vazia ou cheia (dentre outros testes);
\item Temos que inicializar a lista e garantir a ordem de seus elementos.
\end{itemize}
\end{itemize}
\end{frame}

\begin{frame}[fragile]{Modelagem da Lista de Vetor}
\lstset{language=C++,
          keywordstyle=\color{blue}\ttfamily,
          stringstyle=\color{red}\ttfamily,
          commentstyle=\color{green}\ttfamily,
          breaklines=true,
          basicstyle=\ttfamily\footnotesize
          }

\begin{itemize}
\item Inicializar ou limpar:
\begin{itemize}
\item Lista();
\item Lista(int tam);
\item limpaLista();
\item ~Lista();
\end{itemize}
\item Testar se a lista está vazia ou cheia e outros testes:
\begin{itemize}
\item bool pilhaCheia();
\item bool pilhaVazia();
\item int posicao(dado)
\item bool contem(dado)
\end{itemize}
\end{itemize}
\end{frame}

\begin{frame}[fragile]{Modelagem da Lista de Vetor}
\lstset{language=C++,
          keywordstyle=\color{blue}\ttfamily,
          stringstyle=\color{red}\ttfamily,
          commentstyle=\color{green}\ttfamily,
          breaklines=true,
          basicstyle=\ttfamily\footnotesize
          }

\begin{itemize}
\item Colocar e retirar dados da lista:
\begin{itemize}
\item adiciona(T dado);
\item adicionaNoInicio(T dado);
\item adicionaNaPosicao(T dado, int posicao);
\item adicionaEmOrdem(T dado);
\item T retira();
\item T retiraDoInicio();
\item T retiraDaPosicao(int posicao);
\item T retiraEspecifico(dado);
\end{itemize}
\end{itemize}
\end{frame}

\begin{frame}[fragile]{Método \textit{Lista()}}
\lstset{language=C++,
          keywordstyle=\color{blue}\ttfamily,
          stringstyle=\color{red}\ttfamily,
          commentstyle=\color{green}\ttfamily,
          breaklines=true,
          basicstyle=\ttfamily\footnotesize
          }
\begin{lstlisting}
Lista()
inicio
  _dados = new T[_tam];
  _ultimo <- -1;
fim;
\end{lstlisting}
\end{frame}

\begin{frame}[fragile]{Método \textit{Lista(int tam)}}
\lstset{language=C++,
          keywordstyle=\color{blue}\ttfamily,
          stringstyle=\color{red}\ttfamily,
          commentstyle=\color{green}\ttfamily,
          breaklines=true,
          basicstyle=\ttfamily\footnotesize
          }
\begin{lstlisting}
Lista(int tam)
inicio
  _tam = tam;
  _dados = new T[tam];
  ultimo <- -1;
fim;
\end{lstlisting}
\end{frame}

\begin{frame}[fragile]{Método \textit{limpaLista()}}
\lstset{language=C++,
          keywordstyle=\color{blue}\ttfamily,
          stringstyle=\color{red}\ttfamily,
          commentstyle=\color{green}\ttfamily,
          breaklines=true,
          basicstyle=\ttfamily\footnotesize
          }
\begin{lstlisting}
void limpaLista()
inicio
  ultimo <- -1;
fim;
\end{lstlisting}
\end{frame}

\begin{frame}[fragile]{Método \textit{~Lista()}}
\lstset{language=C++,
          keywordstyle=\color{blue}\ttfamily,
          stringstyle=\color{red}\ttfamily,
          commentstyle=\color{green}\ttfamily,
          breaklines=true,
          basicstyle=\ttfamily\footnotesize
          }
\begin{lstlisting}
void ~Lista()
inicio
  delete _dados;
fim;
\end{lstlisting}
\begin{block}{Observação:}
Este método é desnecessário. Ele só está aqui por completude da interface. Mais tarde veremos que, utilizando alocação dinâmica de memória, possuir um destrutor para  Lista é muito importante.
\end{block}
\end{frame}


\begin{frame}[fragile]{Método \textit{listaCheia()}}
\lstset{language=C++,
          keywordstyle=\color{blue}\ttfamily,
          stringstyle=\color{red}\ttfamily,
          commentstyle=\color{green}\ttfamily,
          breaklines=true,
          basicstyle=\ttfamily\footnotesize
          }
\begin{lstlisting}
bool listaCheia()
inicio
 SE (_ultimo = _tam - 1) ENTAO
  RETORNE(Verdadeiro);
 SENAO
  RETORNE(Falso);
fim
\end{lstlisting}
\end{frame}

\begin{frame}[fragile]{Método \textit{listaVazia()}}
\lstset{language=C++,
          keywordstyle=\color{blue}\ttfamily,
          stringstyle=\color{red}\ttfamily,
          commentstyle=\color{green}\ttfamily,
          breaklines=true,
          basicstyle=\ttfamily\footnotesize
          }
\begin{lstlisting}
bool listaVazia()
inicio
SE (_ultimo = -1) ENTAO
 RETORNE(Verdadeiro)
SENAO
 RETORNE(Falso);
fim;
\end{lstlisting}
\end{frame}

\begin{frame}[fragile]{Método \textit{adiciona(T dado)}}
\begin{columns}
\column{.5\textwidth}
\begin{itemize}
\item Testamos se há espaço;
\item Incrementamos o último;
\item Adicionamos o novo dado na posićão último.
\end{itemize}
\column{.5\textwidth}
\begin{center}
\begin{tabular}{| p{.5cm} |p{.5cm}c }
\cline{1-1}
 \cellcolor{SpringGreen}  & &\\ \cline{1-1}
\cellcolor{GreenYellow}  & &\\ \cline{1-1}
\cellcolor{SpringGreen}  12& &\\ \cline{1-1}
\cellcolor{GreenYellow}  4 & &\\ \cline{1-1}
\cellcolor{SpringGreen} 55 & &\\ \cline{1-2}
\cellcolor{GreenYellow} 20 & \multicolumn{1}{|p{.5cm}|} {3} & \\ \cline{1-2}
\end{tabular}
\end{center}
\end{columns}
\end{frame}


\begin{frame}[fragile]{Método \textit{adiciona(T dado)}}
\begin{columns}
\column{.5\textwidth}
\begin{itemize}
\item Testamos se há espaço;
\item Incrementamos o último;
\item Adicionamos o novo dado na posićão último.
\end{itemize}
\column{.5\textwidth}
\begin{center}
\begin{tabular}{| p{.5cm} |p{.5cm}c }
\cline{1-1}
 \cellcolor{SpringGreen}  & &\\ \cline{1-1}
\cellcolor{GreenYellow}  & &\\ \cline{1-1}
\cellcolor{SpringGreen}  12& &\\ \cline{1-1}
\cellcolor{GreenYellow}  4 & &\\ \cline{1-1}
\cellcolor{SpringGreen} 55 & &\\ \cline{1-2}
\cellcolor{GreenYellow} 20 & \multicolumn{1}{|p{.5cm}|} {4} & \\ \cline{1-2}
\end{tabular}
\end{center}
\end{columns}
\end{frame}


\begin{frame}[fragile]{Método \textit{adiciona(T dado)}}
\begin{columns}
\column{.5\textwidth}
\begin{itemize}
\item Testamos se há espaço;
\item Incrementamos o último;
\item Adicionamos o novo dado na posićão último.
\end{itemize}
\column{.5\textwidth}
\begin{center}
\begin{tabular}{| p{.5cm} |p{.5cm}c }
\cline{1-1}
 \cellcolor{SpringGreen}  & &\\ \cline{1-1}
\cellcolor{GreenYellow} 89 & &\\ \cline{1-1}
\cellcolor{SpringGreen}  12& &\\ \cline{1-1}
\cellcolor{GreenYellow}  4 & &\\ \cline{1-1}
\cellcolor{SpringGreen} 55 & &\\ \cline{1-2}
\cellcolor{GreenYellow} 20 & \multicolumn{1}{|p{.5cm}|} {4} & \\ \cline{1-2}
\end{tabular}
\end{center}
\end{columns}
\end{frame}

\begin{frame}[fragile]{Método \textit{adiciona(T dado)}}
\lstset{language=C++,
          keywordstyle=\color{blue}\ttfamily,
          stringstyle=\color{red}\ttfamily,
          commentstyle=\color{green}\ttfamily,
          breaklines=true,
          basicstyle=\ttfamily\footnotesize
          }
\begin{lstlisting}
adiciona(T dado)
inicio
  SE (listaCheia) ENTAO
   THROW(ERROLISTACHEIA)
  SENAO
   _ultimo <- _ultimo + 1;
   _dados[_ultimo] <- dado;
  FIM SE
fim;
\end{lstlisting}
\end{frame}

\begin{frame}[fragile]{Método \textit{T retira()}}
\begin{columns}
\column{.5\textwidth}
\begin{itemize}
\item Testamos se há elementos;
\item Decrementamos o último;
\item Devolvemos o último elemento.
\end{itemize}
\column{.5\textwidth}
\begin{center}
\begin{tabular}{| p{.5cm} |p{.5cm}c }
\cline{1-1}
 \cellcolor{SpringGreen}  & &\\ \cline{1-1}
\cellcolor{GreenYellow} 89 & &\\ \cline{1-1}
\cellcolor{SpringGreen}  12& &\\ \cline{1-1}
\cellcolor{GreenYellow}  4 & &\\ \cline{1-1}
\cellcolor{SpringGreen} 55 & &\\ \cline{1-2}
\cellcolor{GreenYellow} 20 & \multicolumn{1}{|p{.5cm}|} {4} & \\ \cline{1-2}
\end{tabular}
\end{center}
\end{columns}
\end{frame}

\begin{frame}[fragile]{Método \textit{T retira()}}
\begin{columns}
\column{.5\textwidth}
\begin{itemize}
\item Testamos se há elementos;
\item Decrementamos o último;
\item Devolvemos o último elemento.
\end{itemize}
\column{.5\textwidth}
\begin{center}
\begin{tabular}{| p{.5cm} |p{.5cm}c }
\cline{1-1}
 \cellcolor{SpringGreen}  & &\\ \cline{1-1}
\cellcolor{GreenYellow} 89 & &\\ \cline{1-1}
\cellcolor{SpringGreen}  12& &\\ \cline{1-1}
\cellcolor{GreenYellow}  4 & &\\ \cline{1-1}
\cellcolor{SpringGreen} 55 & &\\ \cline{1-2}
\cellcolor{GreenYellow} 20 & \multicolumn{1}{|p{.5cm}|} {3} & \\ \cline{1-2}
\end{tabular}
\end{center}
\end{columns}
\end{frame}

\begin{frame}[fragile]{Método \textit{T retira()}}
\begin{columns}
\column{.5\textwidth}
\begin{itemize}
\item Testamos se há elementos;
\item Decrementamos o último;
\item Devolvemos o último elemento.
\end{itemize}
\column{.5\textwidth}
\begin{center}
\begin{tabular}{| p{.5cm} |p{.5cm}c }
\cline{1-1}
 \cellcolor{SpringGreen}  & &\\ \cline{1-1}
\cellcolor{GreenYellow}  & &\\ \cline{1-1}
\cellcolor{SpringGreen}  12& &\\ \cline{1-1}
\cellcolor{GreenYellow}  4 & &\\ \cline{1-1}
\cellcolor{SpringGreen} 55 & &\\ \cline{1-2}
\cellcolor{GreenYellow} 20 & \multicolumn{1}{|p{.5cm}|} {3} & \\ \cline{1-2}
\end{tabular}
\begin{tabular}{| p{.5cm} | }
\hline
 \cellcolor{Mahogany} {89} \\ \hline
\end{tabular}
\end{center}
\end{columns}
\end{frame}

\begin{frame}[fragile]{Método \textit{T retira()}}
\lstset{language=C++,
          keywordstyle=\color{blue}\ttfamily,
          stringstyle=\color{red}\ttfamily,
          commentstyle=\color{green}\ttfamily,
          breaklines=true,
          basicstyle=\ttfamily\footnotesize
          }
\begin{lstlisting}
T retira()
inicio
 SE (listaVazia) ENTAO
  THROW(ERROLISTAVAZIA)
 SENAO
  _ultimo <- _ultimo - 1;
  RETORNE(_dados[_ultimo + 1]);
 FIM SE
fim;
\end{lstlisting}
\end{frame}

\begin{frame}[fragile]{Método \textit{adicionaNoInicio(T dado)}}
\begin{columns}
\column{.5\textwidth}
\begin{itemize}
\item Testamos se há espaço;
\item Incrementamos o último;
\item Empurramos tudo para trás;
\item Adicionamos o novo dado na primeira posição.
\end{itemize}
\column{.5\textwidth}
\begin{center}
\begin{tabular}{| p{.5cm} |p{.5cm}c }
\cline{1-1}
 \cellcolor{SpringGreen}  & &\\ \cline{1-1}
\cellcolor{GreenYellow} 89 & &\\ \cline{1-1}
\cellcolor{SpringGreen}  12& &\\ \cline{1-1}
\cellcolor{GreenYellow}  4 & &\\ \cline{1-1}
\cellcolor{SpringGreen} 55 & &\\ \cline{1-2}
\cellcolor{GreenYellow} 20 & \multicolumn{1}{|p{.5cm}|} {4} & \\ \cline{1-2}
\end{tabular}
\end{center}
\end{columns}
\end{frame}

\begin{frame}[fragile]{Método \textit{adicionaNoInicio(T dado)}}
\begin{columns}
\column{.5\textwidth}
\begin{itemize}
\item Testamos se há espaço;
\item Incrementamos o último;
\item Empurramos tudo para trás;
\item Adicionamos o novo dado na primeira posição.
\end{itemize}
\column{.5\textwidth}
\begin{center}
\begin{tabular}{| p{.5cm} |p{.5cm}c }
\cline{1-1}
 \cellcolor{SpringGreen}  & &\\ \cline{1-1}
\cellcolor{GreenYellow} 89 & &\\ \cline{1-1}
\cellcolor{SpringGreen}  12& &\\ \cline{1-1}
\cellcolor{GreenYellow}  4 & &\\ \cline{1-1}
\cellcolor{SpringGreen} 55 & &\\ \cline{1-2}
\cellcolor{GreenYellow} 20 & \multicolumn{1}{|p{.5cm}|} {5} & \\ \cline{1-2}
\end{tabular}
\end{center}
\end{columns}
\end{frame}

\begin{frame}[fragile]{Método \textit{adicionaNoInicio(T dado)}}
\begin{columns}
\column{.5\textwidth}
\begin{itemize}
\item Testamos se há espaço;
\item Incrementamos o último;
\item Empurramos tudo para trás;
\item Adicionamos o novo dado na primeira posição.
\end{itemize}
\column{.5\textwidth}
\begin{center}
\begin{tabular}{| p{.5cm} |p{.5cm}c }
\cline{1-1}
 \cellcolor{SpringGreen} 89 & &\\ \cline{1-1}
\cellcolor{GreenYellow}  & &\\ \cline{1-1}
\cellcolor{SpringGreen}  12& &\\ \cline{1-1}
\cellcolor{GreenYellow}  4 & &\\ \cline{1-1}
\cellcolor{SpringGreen} 55 & &\\ \cline{1-2}
\cellcolor{GreenYellow} 20 & \multicolumn{1}{|p{.5cm}|} {5} & \\ \cline{1-2}
\end{tabular}
\end{center}
\end{columns}
\end{frame}

\begin{frame}[fragile]{Método \textit{adicionaNoInicio(T dado)}}
\begin{columns}
\column{.5\textwidth}
\begin{itemize}
\item Testamos se há espaço;
\item Incrementamos o último;
\item Empurramos tudo para trás;
\item Adicionamos o novo dado na primeira posição.
\end{itemize}
\column{.5\textwidth}
\begin{center}
\begin{tabular}{| p{.5cm} |p{.5cm}c }
\cline{1-1}
 \cellcolor{SpringGreen} 89 & &\\ \cline{1-1}
\cellcolor{GreenYellow}  12 & &\\ \cline{1-1}
\cellcolor{SpringGreen}  & &\\ \cline{1-1}
\cellcolor{GreenYellow}  4 & &\\ \cline{1-1}
\cellcolor{SpringGreen} 55 & &\\ \cline{1-2}
\cellcolor{GreenYellow} 20 & \multicolumn{1}{|p{.5cm}|} {5} & \\ \cline{1-2}
\end{tabular}
\end{center}
\end{columns}
\end{frame}

\begin{frame}[fragile]{Método \textit{adicionaNoInicio(T dado)}}
\begin{columns}
\column{.5\textwidth}
\begin{itemize}
\item Testamos se há espaço;
\item Incrementamos o último;
\item Empurramos tudo para trás;
\item Adicionamos o novo dado na primeira posição.
\end{itemize}
\column{.5\textwidth}
\begin{center}
\begin{tabular}{| p{.5cm} |p{.5cm}c }
\cline{1-1}
 \cellcolor{SpringGreen} 89 & &\\ \cline{1-1}
\cellcolor{GreenYellow}  12 & &\\ \cline{1-1}
\cellcolor{SpringGreen}  4 & &\\ \cline{1-1}
\cellcolor{GreenYellow}   & &\\ \cline{1-1}
\cellcolor{SpringGreen} 55 & &\\ \cline{1-2}
\cellcolor{GreenYellow} 20 & \multicolumn{1}{|p{.5cm}|} {5} & \\ \cline{1-2}
\end{tabular}
\end{center}
\end{columns}
\end{frame}

\begin{frame}[fragile]{Método \textit{adicionaNoInicio(T dado)}}
\begin{columns}
\column{.5\textwidth}
\begin{itemize}
\item Testamos se há espaço;
\item Incrementamos o último;
\item Empurramos tudo para trás;
\item Adicionamos o novo dado na primeira posição.
\end{itemize}
\column{.5\textwidth}
\begin{center}
\begin{tabular}{| p{.5cm} |p{.5cm}c }
\cline{1-1}
 \cellcolor{SpringGreen} 89 & &\\ \cline{1-1}
\cellcolor{GreenYellow}  12 & &\\ \cline{1-1}
\cellcolor{SpringGreen}  4 & &\\ \cline{1-1}
\cellcolor{GreenYellow}  55 & &\\ \cline{1-1}
\cellcolor{SpringGreen}  & &\\ \cline{1-2}
\cellcolor{GreenYellow} 20 & \multicolumn{1}{|p{.5cm}|} {5} & \\ \cline{1-2}
\end{tabular}
\end{center}
\end{columns}
\end{frame}

\begin{frame}[fragile]{Método \textit{adicionaNoInicio(T dado)}}
\begin{columns}
\column{.5\textwidth}
\begin{itemize}
\item Testamos se há espaço;
\item Incrementamos o último;
\item Empurramos tudo para trás;
\item Adicionamos o novo dado na primeira posição.
\end{itemize}
\column{.5\textwidth}
\begin{center}
\begin{tabular}{| p{.5cm} |p{.5cm}c }
\cline{1-1}
 \cellcolor{SpringGreen} 89 & &\\ \cline{1-1}
\cellcolor{GreenYellow}  12 & &\\ \cline{1-1}
\cellcolor{SpringGreen}  4 & &\\ \cline{1-1}
\cellcolor{GreenYellow}  55 & &\\ \cline{1-1}
\cellcolor{SpringGreen} 20 & &\\ \cline{1-2}
\cellcolor{GreenYellow}  & \multicolumn{1}{|p{.5cm}|} {5} & \\ \cline{1-2}
\end{tabular}
\end{center}
\end{columns}
\end{frame}

\begin{frame}[fragile]{Método \textit{adicionaNoInicio(T dado)}}
\begin{columns}
\column{.5\textwidth}
\begin{itemize}
\item Testamos se há espaço;
\item Incrementamos o último;
\item Empurramos tudo para trás;
\item Adicionamos o novo dado na primeira posição.
\end{itemize}
\column{.5\textwidth}
\begin{center}
\begin{tabular}{| p{.5cm} |p{.5cm}c }
\cline{1-1}
 \cellcolor{SpringGreen} 89 & &\\ \cline{1-1}
\cellcolor{GreenYellow}  12 & &\\ \cline{1-1}
\cellcolor{SpringGreen}  4 & &\\ \cline{1-1}
\cellcolor{GreenYellow}  55 & &\\ \cline{1-1}
\cellcolor{SpringGreen} 20 & &\\ \cline{1-2}
\cellcolor{GreenYellow} 42 & \multicolumn{1}{|p{.5cm}|} {5} & \\ \cline{1-2}
\end{tabular}
\end{center}
\end{columns}
\end{frame}

\begin{frame}[fragile]{Método \textit{adicionaNoInicio(T dado)}}
\lstset{language=C++,
          keywordstyle=\color{blue}\ttfamily,
          stringstyle=\color{red}\ttfamily,
          commentstyle=\color{green}\ttfamily,
          breaklines=true,
          escapeinside={(*@}{@*)},
          basicstyle=\ttfamily\footnotesize
          }
\begin{lstlisting}
adicionaNoInicio(T dado)
 int posicao; //Var auxiliar para "caminhar"
 inicio
  SE (listaCheia) ENTAO
   THROW(ERROLISTACHEIA)
  SENAO
   _ultimo <- _ultimo + 1;
   posicao <- _ultimo;
   ENQUANTO (posicao > 0) FACA
    //Empurrar tudo para tras
    _dados[posicao] <- _dados[posicao - 1];
    posicao <- posicao - 1;
   FIM ENQUANTO
   _dados[0] <- dado;
  FIM SE
fim;
\end{lstlisting}
\end{frame}

\begin{frame}[fragile]{Método \textit{adicionaNoInicio(T dado)}}
\lstset{language=C++,
          keywordstyle=\color{blue}\ttfamily,
          stringstyle=\color{red}\ttfamily,
          commentstyle=\color{green}\ttfamily,
          breaklines=true,
          escapeinside={(*@}{@*)},
          basicstyle=\ttfamily\footnotesize
          }
\begin{lstlisting}
adicionaNoInicio(T dado)
 int posicao; //Var auxiliar para "caminhar"
 inicio
   (*@  \textcolor{red}{SE (listaCheia) ENTAO}  @*)
    (*@  \textcolor{red}{THROW(ERROLISTACHEIA)}  @*)
   (*@  \textcolor{red}{SENAO}  @*)
   _ultimo <- _ultimo + 1;
   posicao <- _ultimo;
   ENQUANTO (posicao > 0) FACA
    //Empurrar tudo para tras
    _dados[posicao] <- _dados[posicao - 1];
    posicao <- posicao - 1;
   FIM ENQUANTO
   _dados[0] <- dado;
  FIM SE
fim;
\end{lstlisting}
\end{frame}

\begin{frame}[fragile]{Método \textit{adicionaNoInicio(T dado)}}
\lstset{language=C++,
          keywordstyle=\color{blue}\ttfamily,
          stringstyle=\color{red}\ttfamily,
          commentstyle=\color{green}\ttfamily,
          breaklines=true,
          escapeinside={(*@}{@*)},
          basicstyle=\ttfamily\footnotesize
          }
\begin{lstlisting}
adicionaNoInicio(T dado)
 int posicao; //Var auxiliar para "caminhar"
 inicio
   SE (listaCheia) ENTAO)
    THROW(ERROLISTACHEIA)
   SENAO
   (*@  \textcolor{red}{\_ultimo <- \_ultimo + 1;}  @*)
   (*@  \textcolor{red}{posicao <- \_ultimo;}  @*)
   (*@  \textcolor{red}{ENQUANTO (posicao > 0) FACA}  @*)
    (*@  \textcolor{red}{//Empurrar tudo para tras}  @*)
    (*@  \textcolor{red}{\_dados[posicao] <- \_dados[posicao - 1];}  @*)
    (*@  \textcolor{red}{posicao <- posicao - 1;}  @*)
   (*@  \textcolor{red}{FIM ENQUANTO}  @*)
   _dados[0] <- dado;
  FIM SE
fim;
\end{lstlisting}
\end{frame}

\begin{frame}[fragile]{Método \textit{adicionaNoInicio(T dado)}}
\lstset{language=C++,
          keywordstyle=\color{blue}\ttfamily,
          stringstyle=\color{red}\ttfamily,
          commentstyle=\color{green}\ttfamily,
          breaklines=true,
          escapeinside={(*@}{@*)},
          basicstyle=\ttfamily\footnotesize
          }
\begin{lstlisting}
adicionaNoInicio(T dado)
 int posicao; //Var auxiliar para "caminhar"
 inicio
   SE (listaCheia) ENTAO
    THROW(ERROLISTACHEIA)
   SENAO
   _ultimo <- _ultimo + 1;
   posicao <- _ultimo;
   ENQUANTO (posicao > 0) FACA
    //Empurrar tudo para tras
    _dados[posicao] <- _dados[posicao - 1];
    posicao <- posicao - 1;
   FIM ENQUANTO
   (*@  \textcolor{red}{\_dados[0] <- dado;}  @*)
  FIM SE
fim;
\end{lstlisting}
\end{frame}

\begin{frame}[fragile]{Método \textit{T retiraDoInicio()}}
\begin{columns}
\column{.5\textwidth}
\begin{itemize}
\item Testamos se há elementos;
\item Decrementamos o último;
\item Salvamos o primeiro elemento;
\item Empurramos tudo para a frente.
\end{itemize}
\column{.5\textwidth}
\begin{center}
\begin{tabular}{| p{.5cm} |p{.5cm}c }
\cline{1-1}
 \cellcolor{SpringGreen} 89 & &\\ \cline{1-1}
\cellcolor{GreenYellow}  12 & &\\ \cline{1-1}
\cellcolor{SpringGreen}  4 & &\\ \cline{1-1}
\cellcolor{GreenYellow}  55 & &\\ \cline{1-1}
\cellcolor{SpringGreen} 20 & &\\ \cline{1-2}
\cellcolor{GreenYellow} 42 & \multicolumn{1}{|p{.5cm}|} {5} & \\ \cline{1-2}
\end{tabular}
\end{center}
\end{columns}
\end{frame}

\begin{frame}[fragile]{Método \textit{T retiraDoInicio()}}
\begin{columns}
\column{.5\textwidth}
\begin{itemize}
\item Testamos se há elementos;
\item Decrementamos o último;
\item Salvamos o primeiro elemento;
\item Empurramos tudo para a frente.
\end{itemize}
\column{.5\textwidth}
\begin{center}
\begin{tabular}{| p{.5cm} |p{.5cm}c }
\cline{1-1}
 \cellcolor{SpringGreen} 89 & &\\ \cline{1-1}
\cellcolor{GreenYellow}  12 & &\\ \cline{1-1}
\cellcolor{SpringGreen}  4 & &\\ \cline{1-1}
\cellcolor{GreenYellow}  55 & &\\ \cline{1-1}
\cellcolor{SpringGreen} 20 & &\\ \cline{1-2}
\cellcolor{GreenYellow} 42 & \multicolumn{1}{|p{.5cm}|} {4} & \\ \cline{1-2}
\end{tabular}
\end{center}
\end{columns}
\end{frame}

\begin{frame}[fragile]{Método \textit{T retiraDoInicio()}}
\begin{columns}
\column{.5\textwidth}
\begin{itemize}
\item Testamos se há elementos;
\item Decrementamos o último;
\item Salvamos o primeiro elemento;
\item Empurramos tudo para a frente.
\end{itemize}
\column{.5\textwidth}
\begin{center}
\begin{tabular}{| p{.5cm} |p{.5cm}c }
\cline{1-1}
 \cellcolor{SpringGreen} 89 & &\\ \cline{1-1}
\cellcolor{GreenYellow}  12 & &\\ \cline{1-1}
\cellcolor{SpringGreen}  4 & &\\ \cline{1-1}
\cellcolor{GreenYellow}  55 & &\\ \cline{1-1}
\cellcolor{SpringGreen} 20 & &\\ \cline{1-2}
\cellcolor{GreenYellow} 42 & \multicolumn{1}{|p{.5cm}|} {4} & \\ \cline{1-2}
\end{tabular}
\begin{tabular}{| p{.5cm} | }
\hline
 \cellcolor{Mahogany} {42} \\ \hline
\end{tabular}
\end{center}
\end{columns}
\end{frame}

\begin{frame}[fragile]{Método \textit{T retiraDoInicio()}}
\begin{columns}
\column{.5\textwidth}
\begin{itemize}
\item Testamos se há elementos;
\item Decrementamos o último;
\item Salvamos o primeiro elemento;
\item Empurramos tudo para a frente.
\end{itemize}
\column{.5\textwidth}
\begin{center}
\begin{tabular}{| p{.5cm} |p{.5cm}c }
\cline{1-1}
 \cellcolor{SpringGreen} 89 & &\\ \cline{1-1}
\cellcolor{GreenYellow}  12 & &\\ \cline{1-1}
\cellcolor{SpringGreen}  4 & &\\ \cline{1-1}
\cellcolor{GreenYellow}  55 & &\\ \cline{1-1}
\cellcolor{SpringGreen} 20 & &\\ \cline{1-2}
\cellcolor{GreenYellow} 20 & \multicolumn{1}{|p{.5cm}|} {4} & \\ \cline{1-2}
\end{tabular}
\begin{tabular}{| p{.5cm} | }
\hline
 \cellcolor{Mahogany} {42} \\ \hline
\end{tabular}
\end{center}
\end{columns}
\end{frame}

\begin{frame}[fragile]{Método \textit{T retiraDoInicio()}}
\begin{columns}
\column{.5\textwidth}
\begin{itemize}
\item Testamos se há elementos;
\item Decrementamos o último;
\item Salvamos o primeiro elemento;
\item Empurramos tudo para a frente.
\end{itemize}
\column{.5\textwidth}
\begin{center}
\begin{tabular}{| p{.5cm} |p{.5cm}c }
\cline{1-1}
 \cellcolor{SpringGreen} 89 & &\\ \cline{1-1}
\cellcolor{GreenYellow}  12 & &\\ \cline{1-1}
\cellcolor{SpringGreen}  4 & &\\ \cline{1-1}
\cellcolor{GreenYellow}  55 & &\\ \cline{1-1}
\cellcolor{SpringGreen} 55 & &\\ \cline{1-2}
\cellcolor{GreenYellow} 20 & \multicolumn{1}{|p{.5cm}|} {4} & \\ \cline{1-2}
\end{tabular}
\begin{tabular}{| p{.5cm} | }
\hline
 \cellcolor{Mahogany} {42} \\ \hline
\end{tabular}
\end{center}
\end{columns}
\end{frame}

\begin{frame}[fragile]{Método \textit{T retiraDoInicio()}}
\begin{columns}
\column{.5\textwidth}
\begin{itemize}
\item Testamos se há elementos;
\item Decrementamos o último;
\item Salvamos o primeiro elemento;
\item Empurramos tudo para a frente.
\end{itemize}
\column{.5\textwidth}
\begin{center}
\begin{tabular}{| p{.5cm} |p{.5cm}c }
\cline{1-1}
 \cellcolor{SpringGreen} 89 & &\\ \cline{1-1}
\cellcolor{GreenYellow}  12 & &\\ \cline{1-1}
\cellcolor{SpringGreen}  4 & &\\ \cline{1-1}
\cellcolor{GreenYellow}  4 & &\\ \cline{1-1}
\cellcolor{SpringGreen} 55 & &\\ \cline{1-2}
\cellcolor{GreenYellow} 20 & \multicolumn{1}{|p{.5cm}|} {4} & \\ \cline{1-2}
\end{tabular}
\begin{tabular}{| p{.5cm} | }
\hline
 \cellcolor{Mahogany} {42} \\ \hline
\end{tabular}
\end{center}
\end{columns}
\end{frame}

\begin{frame}[fragile]{Método \textit{T retiraDoInicio()}}
\begin{columns}
\column{.5\textwidth}
\begin{itemize}
\item Testamos se há elementos;
\item Decrementamos o último;
\item Salvamos o primeiro elemento;
\item Empurramos tudo para a frente.
\end{itemize}
\column{.5\textwidth}
\begin{center}
\begin{tabular}{| p{.5cm} |p{.5cm}c }
\cline{1-1}
 \cellcolor{SpringGreen} 89 & &\\ \cline{1-1}
\cellcolor{GreenYellow}  12 & &\\ \cline{1-1}
\cellcolor{SpringGreen}  12 & &\\ \cline{1-1}
\cellcolor{GreenYellow}  4 & &\\ \cline{1-1}
\cellcolor{SpringGreen} 55 & &\\ \cline{1-2}
\cellcolor{GreenYellow} 20 & \multicolumn{1}{|p{.5cm}|} {4} & \\ \cline{1-2}
\end{tabular}
\begin{tabular}{| p{.5cm} | }
\hline
 \cellcolor{Mahogany} {42} \\ \hline
\end{tabular}
\end{center}
\end{columns}
\end{frame}

\begin{frame}[fragile]{Método \textit{T retiraDoInicio()}}
\begin{columns}
\column{.5\textwidth}
\begin{itemize}
\item Testamos se há elementos;
\item Decrementamos o último;
\item Salvamos o primeiro elemento;
\item Empurramos tudo para a frente.
\end{itemize}
\column{.5\textwidth}
\begin{center}
\begin{tabular}{| p{.5cm} |p{.5cm}c }
\cline{1-1}
 \cellcolor{SpringGreen} 89 & &\\ \cline{1-1}
\cellcolor{GreenYellow}  89 & &\\ \cline{1-1}
\cellcolor{SpringGreen}  12 & &\\ \cline{1-1}
\cellcolor{GreenYellow}  4 & &\\ \cline{1-1}
\cellcolor{SpringGreen} 55 & &\\ \cline{1-2}
\cellcolor{GreenYellow} 20 & \multicolumn{1}{|p{.5cm}|} {4} & \\ \cline{1-2}
\end{tabular}
\begin{tabular}{| p{.5cm} | }
\hline
 \cellcolor{Mahogany} {42} \\ \hline
\end{tabular}
\end{center}
\end{columns}
\end{frame}

\begin{frame}[fragile]{Método \textit{T retiraDoInicio()}}
\begin{columns}
\column{.5\textwidth}
\begin{itemize}
\item Testamos se há elementos;
\item Decrementamos o último;
\item Salvamos o primeiro elemento;
\item Empurramos tudo para a frente.
\item{\color{red} Não perdemos tempo apagando o segundo "89", pois o \_ultimo determina que a lista acaba antes dele!}
\end{itemize}
\column{.5\textwidth}
\begin{center}
\begin{tabular}{| p{.5cm} |p{.5cm}c }
\cline{1-1}
 \cellcolor{SpringGreen} 89 & &\\ \cline{1-1}
\cellcolor{GreenYellow}  89 & &\\ \cline{1-1}
\cellcolor{SpringGreen}  12 & &\\ \cline{1-1}
\cellcolor{GreenYellow}  4 & &\\ \cline{1-1}
\cellcolor{SpringGreen} 55 & &\\ \cline{1-2}
\cellcolor{GreenYellow} 20 & \multicolumn{1}{|p{.5cm}|} {4} & \\ \cline{1-2}
\end{tabular}
\begin{tabular}{| p{.5cm} | }
\hline
 \cellcolor{Mahogany} {42} \\ \hline
\end{tabular}
\end{center}
\end{columns}
\end{frame}

\begin{frame}[fragile]{Método \textit{T retiraDoInicio()}}
\lstset{language=C++,
          keywordstyle=\color{blue}\ttfamily,
          stringstyle=\color{red}\ttfamily,
          commentstyle=\color{green}\ttfamily,
          breaklines=true,
          escapeinside={(*@}{@*)},
          basicstyle=\ttfamily\footnotesize
          }
\begin{lstlisting}
T retiraDoInicio()
 int posicao; T valor;
inicio
 SE (listaVazia) ENTAO
  THROW(ERROLISTAVAZIA)
 SENAO
  _ultimo <- _ultimo - 1;
  valor <- _dados[0];
  posicao <- 0;
  ENQUANTO (posicao <= _ultimo) FACA
   //Empurrar tudo para a frente
   _dados[posicao] <- _dados[posicao + 1];
   posicao <- posicao + 1;
  FIM ENQUANTO
  RETORNE(valor);
 FIM SE
fim;
(*@  \textcolor{red}{ }  @*)
\end{lstlisting}
\end{frame}

\begin{frame}[fragile]{Método \textit{T retiraDoInicio()}}
\lstset{language=C++,
          keywordstyle=\color{blue}\ttfamily,
          stringstyle=\color{red}\ttfamily,
          commentstyle=\color{green}\ttfamily,
          breaklines=true,
          escapeinside={(*@}{@*)},
          basicstyle=\ttfamily\footnotesize
          }
\begin{lstlisting}
T retiraDoInicio()
 int posicao; T valor;
inicio
 (*@  \textcolor{red}{SE (listaVazia) ENTAO}  @*)
  (*@  \textcolor{red}{THROW(ERROLISTAVAZIA)}  @*)
 (*@  \textcolor{red}{SENAO}  @*)
  _ultimo <- _ultimo - 1;
  valor <- _dados[0];
  posicao <- 0;
  ENQUANTO (posicao <= _ultimo) FACA
   //Empurrar tudo para a frente
   _dados[posicao] <- _dados[posicao + 1];
   posicao <- posicao + 1;
  FIM ENQUANTO
  RETORNE(valor);
 FIM SE
fim;
(*@  \textcolor{red}{ }  @*)
\end{lstlisting}
\end{frame}

\begin{frame}[fragile]{Método \textit{T retiraDoInicio()}}
\lstset{language=C++,
          keywordstyle=\color{blue}\ttfamily,
          stringstyle=\color{red}\ttfamily,
          commentstyle=\color{green}\ttfamily,
          breaklines=true,
          escapeinside={(*@}{@*)},
          basicstyle=\ttfamily\footnotesize
          }
\begin{lstlisting}
T retiraDoInicio()
 int posicao; T valor;
inicio
 SE (listaVazia) ENTAO
  THROW(ERROLISTAVAZIA)
 SENAO
  _ultimo <- _ultimo - 1;
  (*@  \textcolor{red}{valor <- \_dados[0];}  @*)
  posicao <- 0;
  ENQUANTO (posicao <= _ultimo) FACA
   //Empurrar tudo para a frente
   _dados[posicao] <- _dados[_posicao + 1];
   posicao <- posicao + 1;
  FIM ENQUANTO
  RETORNE(valor);
 FIM SE
fim;

\end{lstlisting}
\end{frame}

\begin{frame}[fragile]{Método \textit{T retiraDoInicio()}}
\lstset{language=C++,
          keywordstyle=\color{blue}\ttfamily,
          stringstyle=\color{red}\ttfamily,
          commentstyle=\color{green}\ttfamily,
          breaklines=true,
          escapeinside={(*@}{@*)},
          basicstyle=\ttfamily\footnotesize
          }
\begin{lstlisting}
T retiraDoInicio()
 int posicao; T valor;
inicio
 SE (listaVazia) ENTAO
  THROW(ERROLISTAVAZIA)
 SENAO
  _ultimo <- _ultimo - 1;
  valor <- _dados[0];
  (*@  \textcolor{red}{posicao <- 0;}  @*)
  (*@  \textcolor{red}{ENQUANTO (posicao <= \_ultimo) FACA }  @*)
   (*@  \textcolor{red}{//Empurrar tudo para a frente}  @*)
   (*@  \textcolor{red}{\_dados[posicao] <- \_dados[posicao + 1];}  @*)
   (*@  \textcolor{red}{posicao <- posicao + 1;}  @*)
  (*@  \textcolor{red}{FIM ENQUANTO}  @*)
  RETORNE(valor);
 FIM SE
fim;

\end{lstlisting}
\end{frame}

\begin{frame}[fragile]{Método \textit{T retiraDoInicio()}}
\lstset{language=C++,
          keywordstyle=\color{blue}\ttfamily,
          stringstyle=\color{red}\ttfamily,
          commentstyle=\color{green}\ttfamily,
          breaklines=true,
          escapeinside={(*@}{@*)},
          basicstyle=\ttfamily\footnotesize
          }
\begin{lstlisting}
T retiraDoInicio()
 int posicao; T valor;
inicio
 SE (listaVazia) ENTAO
  THROW(ERROLISTAVAZIA)
 SENAO
  _ultimo <- _ultimo - 1;
  valor <- _dados[0];
  posicao <- 0;
  ENQUANTO (posicao <= _ultimo) FACA
   //Empurrar tudo para a frente
   _dados[_posicao] <- _dados[_posicao + 1];
   _posicao <- posicao + 1;
  FIM ENQUANTO
  (*@  \textcolor{red}{RETORNE(valor);}  @*)
 FIM SE
fim;

\end{lstlisting}
\end{frame}

\begin{frame}[fragile]{Método \textit{adicionaNaPosicao(T dado, int posicao)}}
\begin{columns}
\column{.5\textwidth}
\begin{itemize}
\item Testamos se há espaço e se a posição existe;
\item Incrementamos o último;
\item Empurramos tudo para trás a partir da posição;
\item Adicionamos o novo dado na posição informada.
\end{itemize}
\column{.5\textwidth}
\begin{center}
\begin{tabular}{| p{.5cm} |p{.5cm}c }
\cline{1-1}
 \cellcolor{SpringGreen}  & &\\ \cline{1-1}
\cellcolor{GreenYellow}  89 & &\\ \cline{1-1}
\cellcolor{SpringGreen}  12 & &\\ \cline{1-1}
\cellcolor{GreenYellow}  4 & &\\ \cline{1-1}
\cellcolor{SpringGreen} 55 & &\\ \cline{1-2}
\cellcolor{GreenYellow} 20 & \multicolumn{1}{|p{.5cm}|} {4} & \\ \cline{1-2}
\end{tabular}
\end{center}
\end{columns}
\end{frame}

\begin{frame}[fragile]{Método \textit{adicionaNaPosicao(T dado, int posicao)}}
\begin{columns}
\column{.5\textwidth}
\begin{itemize}
\item Testamos se há espaço e se a posição existe;
\item Incrementamos o último;
\item Empurramos tudo para trás a partir da posição;
\item Adicionamos o novo dado na posição informada.
\end{itemize}
\column{.5\textwidth}
\begin{center}
\begin{tabular}{| p{.5cm} |p{.5cm}c }
\cline{1-1}
 \cellcolor{SpringGreen}  & &\\ \cline{1-1}
\cellcolor{GreenYellow}  89 & &\\ \cline{1-1}
\cellcolor{SpringGreen}  12 & &\\ \cline{1-1}
\cellcolor{GreenYellow}  4 & &\\ \cline{1-1}
\cellcolor{SpringGreen} 55 & &\\ \cline{1-2}
\cellcolor{GreenYellow} 20 & \multicolumn{1}{|p{.5cm}|} {5} & \\ \cline{1-2}
\end{tabular}
\end{center}
\end{columns}
\end{frame}

\begin{frame}[fragile]{Método \textit{adicionaNaPosicao(T dado, int posicao)}}
\begin{columns}
\column{.5\textwidth}
\begin{itemize}
\item Testamos se há espaço e se a posição existe;
\item Incrementamos o último;
\item Empurramos tudo para trás a partir da posição;
\item Adicionamos o novo dado na posição informada.
\end{itemize}
\column{.5\textwidth}
\begin{center}
\begin{tabular}{| p{.5cm} |p{.5cm}c }
\cline{1-1}
 \cellcolor{SpringGreen} 89 & &\\ \cline{1-1}
\cellcolor{GreenYellow}  89 & &\\ \cline{1-1}
\cellcolor{SpringGreen}  12 & &\\ \cline{1-1}
\cellcolor{GreenYellow}  4 & &\\ \cline{1-1}
\cellcolor{SpringGreen} 55 & &\\ \cline{1-2}
\cellcolor{GreenYellow} 20 & \multicolumn{1}{|p{.5cm}|} {5} & \\ \cline{1-2}
\end{tabular}
\end{center}
\end{columns}
\end{frame}

\begin{frame}[fragile]{Método \textit{adicionaNaPosicao(T dado, int posicao)}}
\begin{columns}
\column{.5\textwidth}
\begin{itemize}
\item Testamos se há espaço e se a posição existe;
\item Incrementamos o último;
\item Empurramos tudo para trás a partir da posição;
\item Adicionamos o novo dado na posição informada.
\end{itemize}
\column{.5\textwidth}
\begin{center}
\begin{tabular}{| p{.5cm} |p{.5cm}c }
\cline{1-1}
 \cellcolor{SpringGreen} 89 & &\\ \cline{1-1}
\cellcolor{GreenYellow}  12 & &\\ \cline{1-1}
\cellcolor{SpringGreen}  12 & &\\ \cline{1-1}
\cellcolor{GreenYellow}  4 & &\\ \cline{1-1}
\cellcolor{SpringGreen} 55 & &\\ \cline{1-2}
\cellcolor{GreenYellow} 20 & \multicolumn{1}{|p{.5cm}|} {5} & \\ \cline{1-2}
\end{tabular}
\end{center}
\end{columns}
\end{frame}

\begin{frame}[fragile]{Método \textit{adicionaNaPosicao(T dado, int posicao)}}
\begin{columns}
\column{.5\textwidth}
\begin{itemize}
\item Testamos se há espaço e se a posição existe;
\item Incrementamos o último;
\item Empurramos tudo para trás a partir da posição;
\item Adicionamos o novo dado na posição informada.
\end{itemize}
\column{.5\textwidth}
\begin{center}
\begin{tabular}{| p{.5cm} |p{.5cm}c }
\cline{1-1}
 \cellcolor{SpringGreen} 89 & &\\ \cline{1-1}
\cellcolor{GreenYellow}  12 & &\\ \cline{1-1}
\cellcolor{SpringGreen}  42 & &\\ \cline{1-1}
\cellcolor{GreenYellow}  4 & &\\ \cline{1-1}
\cellcolor{SpringGreen} 55 & &\\ \cline{1-2}
\cellcolor{GreenYellow} 20 & \multicolumn{1}{|p{.5cm}|} {5} & \\ \cline{1-2}
\end{tabular}
\end{center}
\end{columns}
\end{frame}

\begin{frame}[fragile]{Método \textit{adicionaNaPosicao(T dado, int posicao)}}
\lstset{language=C++,
          keywordstyle=\color{blue}\ttfamily,
          stringstyle=\color{red}\ttfamily,
          commentstyle=\color{green}\ttfamily,
          breaklines=true,
          escapeinside={(*@}{@*)},
          basicstyle=\ttfamily\footnotesize
          }
\begin{lstlisting}
adicionaNaPosicao(T dado, int posicao)
 int atual;
inicio
 SE (listaCheia) ENTAO
  THROW(ERROLISTACHEIA)
 SENAO
  SE (posicao > _ultimo+1 OU posicao < 0) ENTAO
   THROW(ERROPOSICAO);
  FIM SE
  _ultimo <- _ultimo+1; atual <- _ultimo;
  ENQUANTO (atual > posicao) FACA
   _dados[atual] <- _dados[atual - 1];
   atual <- atual - 1;
  FIM ENQUANTO
  _dados[posicao] <- dado;
 FIM SE
fim;
\end{lstlisting}
\end{frame}

\begin{frame}[fragile]{Método \textit{adicionaNaPosicao(T dado, int posicao)}}
\lstset{language=C++,
          keywordstyle=\color{blue}\ttfamily,
          stringstyle=\color{red}\ttfamily,
          commentstyle=\color{green}\ttfamily,
          breaklines=true,
          escapeinside={(*@}{@*)},
          basicstyle=\ttfamily\footnotesize
          }
\begin{lstlisting}
adicionaNaPosicao(T dado, int posicao)
 int atual;
inicio
 SE (listaCheia) ENTAO
  THROW(ERROLISTACHEIA)
 SENAO
  (*@  \textcolor{red}{SE (posicao > \_ultimo+1 OU posicao < 0) ENTAO}  @*)
   (*@  \textcolor{red}{THROW(ERROPOSICAO);}  @*)
  (*@  \textcolor{red}{FIM SE}  @*)
  _ultimo <- _ultimo+1; atual <- _ultimo;
  ENQUANTO (atual > posicao) FACA
   _dados[atual] <- _dados[atual - 1];
   atual <- atual - 1;
  FIM ENQUANTO
  _dados[posicao] <- dado;
 FIM SE
fim;
\end{lstlisting}
\end{frame}

\begin{frame}[fragile]{Método \textit{adicionaNaPosicao(T dado, int posicao)}}
\lstset{language=C++,
          keywordstyle=\color{blue}\ttfamily,
          stringstyle=\color{red}\ttfamily,
          commentstyle=\color{green}\ttfamily,
          breaklines=true,
          escapeinside={(*@}{@*)},
          basicstyle=\ttfamily\footnotesize
          }
\begin{lstlisting}
adicionaNaPosicao(T dado, int posicao)
 int atual;
inicio
 SE (listaCheia) ENTAO
  THROW(ERROLISTACHEIA)
 SENAO
  SE (posicao > _ultimo+1 OU posicao < 0) ENTAO
   THROW(ERROPOSICAO);
  FIM SE
  _ultimo <- _ultimo+1; atual <- _ultimo;
  ENQUANTO (atual > posicao) FACA
   _dados[atual] <- _dados[atual - 1];
   atual <- atual - 1;
  FIM ENQUANTO
  (*@  \textcolor{red}{\_dados[posicao] <- dado;}  @*)
 FIM SE
fim;

\end{lstlisting}
\end{frame}

\begin{frame}[fragile]{Revisitando Métodos}
\begin{itemize}
\item adicionaNoInicio(dado) == adicionaNaPosicao(dado,0);
\item adiciona(dado) == adicionaNaPosicao(dado,ultimo+1);
\end{itemize}
\end{frame}





\begin{frame}[fragile]{Trabalho Lista de Vetor}
\lstset{language=C++,
          keywordstyle=\color{blue}\ttfamily,
          stringstyle=\color{red}\ttfamily,
          commentstyle=\color{green}\ttfamily,
          breaklines=true,
          basicstyle=\ttfamily\footnotesize
          }
\begin{itemize}
\item Implemente uma classe Lista todas as operações vistas;
\item Implemente a lista usando Templates;
\item Implemente a lista com um numero de elementos variável definido na instanciação;
\item Use as melhores práticas de orientação a objetos;
\item Documente todas as classes, métodos e atributos;
\item Aplique os testes unitários disponíveis no moodle da disciplina para validar sua estrutura de dados;
\item Entregue até a data definida no moodle.
\end{itemize}
\end{frame}

{
\usebackgroundtemplate{\includegraphics[width=\paperwidth,
height=\paperheight]{../reusable_images/fundo_capa.png}}
\begin{frame}

{\LARGE Perguntas????}

\end{frame}
}


{
\usebackgroundtemplate{\includegraphics[width=\paperwidth,
height=\paperheight]{../reusable_images/fundo_capa.png}}
\begin{frame}
\includegraphics[scale=0.8]{../reusable_images/cc_logo_arge.png}\hspace{0.5cm}
\includegraphics[scale=0.95]{../reusable_images/by.png}

\vspace{1cm}
Este trabalho está licenciado sob uma Licença Creative Commons Atribuição 4.0 Internacional. Para ver uma cópia desta licença, visite http://creativecommons.org/licenses/by/4.0/.

\end{frame}
}
\end{document}
